\section{磁盘完成请求时怎样唤醒 Thread}

异步 I/O 把命令、缓冲区、完成结果和等待者从调用栈中抽出来，交给一个在
IRQ、超时和取消竞争下仍只有唯一结论的请求状态机。可写共享映射也不能只打开
PTE 的 W 位；Kernel 必须先取得脏页
责任，写回时重新建立下一次写通知边界。

\topic{为什么启动链必须轮询，运行期才适合中断}
CPU 从复位向量进入 ROM 时没有 Process、Thread、IDT、PIC 路由、WaitQueue
或 scheduler。ROM 为了读取 Stage 1，只能发出 ATA 命令，再有界检查
BSY、DRQ、ERR 和 DF。Stage 1 也不能等待由尚未装载的 Kernel 处理 IRQ。
因此 polling 是启动依赖图的必然结果，不是尚未学会中断的临时错误。

\begin{lstlisting}
Reset ROM
  -> bounded ATA polling
  -> load Stage 1
  -> bounded ATA polling
  -> load Kernel
  -> IDT + PIC + PIT + scheduler + WaitQueue
  -> IRQ14 completion becomes safe
\end{lstlisting}

运行期继续轮询却会浪费已经存在的并发结构。一个 Thread 等设备时反复读端口，
其他 Ready Thread 无法获得本应可用的 CPU。v1.16 因而保留两种适配器：
ROM、Stage 1 和 early Kernel 以 nIEN 禁止设备中断并有界轮询；调度器就绪后，
运行期请求队列才清 nIEN、开放 IRQ14。两者不能同时拥有 ATA 通道。

早期 PATA 控制器与现代 NVMe 的另一个差别是：legacy primary channel 没有
command tag。IRQ14 本身不会携带“这是第几个请求”。只要硬件里同时存在多个
未完成命令，软件便无法可靠配对。因此软件队列可以有多个 Queued，请求进入
硬件后却只能有一个 Issued。这个单飞约束来自协议，不是数组容量不够。

\topic{从 ATA 引脚语义到 CPU vector 0x2E}
QEMU 的 \code{pc} 机器模拟传统连线。primary ATA controller 把中断接到
slave 8259A 的 IR6；slave 的输出接到 master IR2；master INTR 再经 Local
APIC LINT0 的 ExtINT virtual-wire 模式到达 CPU。

\begin{pdfdiagram}[node distance=8mm and 10mm]
  \node[os hardware] (ata) {ATA primary\\IRQ14};
  \node[os hardware,right=of ata] (slave) {8259A slave\\IR6};
  \node[os hardware,right=of slave] (master) {8259A master\\IR2 cascade};
  \node[os hardware,right=of master] (lapic) {LAPIC LINT0\\ExtINT};
  \node[os state,below=of lapic] (cpu) {CPU vector 0x2E\\IDT interrupt gate};
  \node[os state,left=of cpu] (asm) {NASM IRQ stub\\统一 176 B frame};
  \node[os block,left=of asm] (cpp) {InterruptRuntime\\ATA completion};
  \draw[os arrow] (ata) -- (slave);
  \draw[os arrow] (slave) -- (master);
  \draw[os arrow] (master) -- (lapic);
  \draw[os arrow] (lapic) -- (cpu);
  \draw[os arrow] (cpu) -- (asm);
  \draw[os arrow] (asm) -- (cpp);
\end{pdfdiagram}
\diagramnote{IRQ14 经过 slave 与 master 两层 in-service 状态；EOI 必须先 slave 后 master。}

PIC 初始化后全屏蔽。v1.16 最终掩码是 \code{0xBFF8}：

\begin{table}[htbp]
\centering
\small
\begin{osgridtabularx}{\textwidth}{l l X}
控制器 & mask & 开放输入\\ \hline
master & \code{0xF8 = 11111000b} & IRQ0 PIT、IRQ1 keyboard、IRQ2 cascade\\ \hline
slave & \code{0xBF = 10111111b} & IR6，即系统 IRQ14\\ \hline
\end{osgridtabularx}
\caption{开放 slave IRQ 必须同时开放 master cascade}
\end{table}

若只清 slave bit 6，slave 可以登记 in-service，但 master 仍挡住级联；若只清
master bit 2，级联线可达，却没有允许的 slave 来源。真实 IRQ14 完成后先向
\code{0xA0} 发 EOI，再向 \code{0x20} 发 EOI；相反顺序可能让级联服务状态
停留，下一次 IRQ 看起来像随机丢失。

\topic{ATA 命令块、状态位与控制位}\par
\begin{table}[htbp]
\centering
\small
\begin{osgridtabularx}{\textwidth}{l l l X}
端口 & 宽度 & 名称 & 读写语义\\ \hline
\code{0x1F0} & 16 & DATA & 一个扇区 256 个 word\\ \hline
\code{0x1F1} & 8 & ERROR/FEATURES & 读错误，写 feature\\ \hline
\code{0x1F2} & 8 & SECTOR COUNT & 当前命令写一扇区\\ \hline
\code{0x1F3..0x1F5} & 8 & LBA LOW/MID/HIGH & LBA 的低 24 位\\ \hline
\code{0x1F6} & 8 & DRIVE/HEAD & master、LBA mode 与 LBA[27:24]\\ \hline
\code{0x1F7} & 8 & STATUS/COMMAND & 读状态，写命令\\ \hline
\code{0x3F6} & 8 & ALT STATUS/CONTROL & 读状态不确认 IRQ；写 nIEN/SRST\\ \hline
\end{osgridtabularx}
\caption{primary ATA command block 与 control port}
\end{table}

Status 的决定性位如下：

\begin{table}[htbp]
\centering
\small
\begin{osgridtabularx}{\textwidth}{c l X}
位 & 名称 & v1.16 解释\\ \hline
7 & BSY & 设备仍拥有内部状态，不能发下一命令\\ \hline
6 & DRDY & 设备就绪提示，不单独表示传输完成\\ \hline
5 & DF & device fault，本请求完成为 DeviceError\\ \hline
3 & DRQ & CPU 必须执行一个数据阶段\\ \hline
0 & ERR & 命令错误，读取 ERROR 并完成为 DeviceError\\ \hline
\end{osgridtabularx}
\caption{ATA STATUS 标志不是可任意组合的软件布尔量}
\end{table}

Device Control bit 1 是 nIEN，置一禁止设备 IRQ；bit 2 是 SRST。读取普通
status 可能确认 pending interrupt，读取 alternate status 不确认，所以两种
端口不能仅因返回相似字节便互换。

Read 请求收到 DRQ IRQ 后执行 256 次 16-bit IN；Write 在第一个 DRQ IRQ
执行 256 次 16-bit OUT，随后还要等待设备完成 IRQ；Flush 没有数据阶段。
逐 word 传输不打印日志，否则串口开销会改变设备与调度时序。

\topic{RFLAGS.IF 与页故障 error code}
RFLAGS bit 9 的 IF 控制可屏蔽外部中断。Kernel 在 IDT、PIC、设备对象和请求
存储发布前保持 IF=0，全部准备完成后才 STI。interrupt gate 进入时处理器
清 IF，并把旧 RFLAGS 保存进特权帧；IRETQ 恢复它。C++ handler 不在任意点
重新 STI，避免同一单飞请求被嵌套 IRQ 重入。

文件页的第一次写则依赖 vector 14 的 page-fault error code：

\begin{table}[htbp]
\centering
\small
\begin{osgridtabularx}{\textwidth}{c l l X}
位 & 名称 & 写通知要求 & 排除的情况\\ \hline
0 & P & 1 & not-present 首次装页\\ \hline
1 & W/R & 1 & 读访问\\ \hline
2 & U/S & 1 & Kernel 自身访问\\ \hline
3 & RSVD & 0 & 页表保留位损坏\\ \hline
4 & I/D & 0 & 取指权限 fault\\ \hline
\end{osgridtabularx}
\caption{writable FileShared 只接管 present+write+user 保护 fault}
\end{table}

真正只读文件映射继续失败；FilePrivate 进入 COW；FileShared 才进入
MarkDirty。若把所有写 fault 都直接改 PTE.W，便会同时破坏只读保护、private
隔离和 dirty backpressure。

\topic{请求身份为何不能等于数组下标}
\code{BlockRequestQueue} 使用调用方提供的 64 槽固定存储，不在 IRQ 中申请。
每个请求保存 operation、LBA、Kernel buffer、owner Thread、absolute deadline、
state、result 和 64 位单调 identifier。

\begin{lstlisting}
Unused
  -> Queued
       -> Completed(Cancelled)
       -> Issued
            -> Completed(Succeeded)
            -> Completed(DeviceError)
            -> Completed(TimedOut)
  -> Reap -> Unused
\end{lstlisting}

数组下标会复用。若等待者只保存下标，请求 A 完成后槽位立刻给请求 B，迟醒的
等待者可能把 B 当成 A，这就是 ABA。identifier 与槽位分离，通用队列的消费方
必须在复制冻结结果后显式 Reap，槽位才重新成为 Unused。生产 ATA 适配器在
IRQ/超时路径中完成这次 copy+Reap，再把独立完成记录交给调度器；等待 Thread
不引用已经复用的槽位。

Issued 请求不能普通取消。命令已经进入设备，仅从链表删除不会撤回硬件状态；
稍后 IRQ 会被下一请求误认。Queued 可以取消，因为它尚未进入设备。

容量耗尽时 Submit 返回明确失败。队列不能覆盖最旧请求，也不能无限增长。
64 是本阶段的内存/压力边界，identifier 和累计统计仍使用显式 \code{uint64_t}。

\topic{IRQ 与 PIT 超时怎样只产生一个结果}
IRQ14 和 PIT IRQ0 可能在同一 deadline 附近到达。唯一完成由状态转移保证：

\begin{enumerate}
  \item IRQ14 先把 Issued 变为 Completed，PIT 随后看不到可超时的 Issued；
  \item PIT 先提交 TimedOut，迟到 IRQ 不能再次 Complete 同一请求；
  \item 重复解析增加诊断统计，但不能覆盖首次结果。
\end{enumerate}

超时不仅是返回值。ATA 可能仍 BSY、等待 data word 或准备产生迟到 IRQ。
ResolveTimeout 获胜后执行 software reset：

\begin{lstlisting}
Complete(TimedOut)
  -> SRST = 1
  -> required delay reads
  -> SRST = 0
  -> bounded wait for BSY clear
  -> only then IssueNext
\end{lstlisting}

若只唤醒 Thread 就发下一命令，两个协议阶段会交叉，队列在软件上正确而设备
状态已经损坏。若 reset 的有界 BSY 检查也失败，controller 被永久冻结；
队列余项逐个以 DeviceError 完成，后续 Submit 直接失败，不再向未知状态的
硬件发命令。这样等待者能收敛，设备所有权也不会被伪装成已经恢复。

\topic{BlockIo 等待与调度证据}
显式 sync 在写回 VFS 后提交异步 Flush。调用 Thread 以请求 identifier 进入
\code{WaitCondition::BlockIo}。IRQ handler 只把 owner 从 Blocked 变为 Ready，
不在任意 IRQ 调用栈上直接切换；现有返回用户态前调度边界稍后选择它。

\begin{pdfdiagram}[node distance=8mm and 10mm]
  \node[os state] (a) {Thread A\\Submit Flush};
  \node[os state,right=of a] (blocked) {A: Blocked\\BlockIo};
  \node[os block,right=of blocked] (b) {Thread B\\继续前进};
  \node[os hardware,below=of b] (irq) {ATA IRQ14};
  \node[os state,left=of irq] (ready) {copy + Reap\\匹配 owner/id};
  \node[os state,left=of ready] (reap) {A: Ready\\结果写入现场};
  \draw[os arrow] (a) -- (blocked);
  \draw[os arrow] (blocked) -- (b);
  \draw[os arrow] (b) -- (irq);
  \draw[os arrow] (irq) -- (ready);
  \draw[os arrow] (ready) -- (reap);
\end{pdfdiagram}
\diagramnote{OTHER\_THREAD\_PROGRESS=1 证明等待期间发生了真实前进，而不只证明最终返回。}

普通 signal 不把 BlockIo 当成字符输入等待唤醒。若一个系统调用需要可中断 I/O，
必须另行冻结取消、设备撤回和 restart ABI；v1.16 不伪造该语义。

完成路径同时核对 owner Thread 槽和 64 位 identifier。Thread 已退出或槽位已
复用时，迟到完成成为 abandoned 诊断，不会唤醒后来占用同一槽的新 Thread。
当前 rootfs 扇区 Read/Write 仍走有界同步 PIO；本阶段接入 BlockIo 的真实生产
消费者是显式 sync 的最终 ATA FLUSH。异步 Read/Write 状态机是后续迁移
journal 和块消费者的基础，不能据此宣称全部文件 I/O 已异步化。

\topic{为什么 clean/dirty 两态仍然不够}\par
\begin{table}[htbp]
\centering
\small
\begin{osgridtabularx}{\textwidth}{l X l}
状态 & 责任 & 零引用时能否 clean LRU 淘汰\\ \hline
Empty & 没有 frame 或文件身份 & 不适用\\ \hline
Clean & frame 与文件一致，无写回债务 & 可以\\ \hline
Dirty & 内存比文件新 & 不可以\\ \hline
Writeback & writer 正在提交这个快照 & 不可以\\ \hline
Error & 上次 writer 失败，债务仍存在 & 不可以\\ \hline
\end{osgridtabularx}
\caption{Error 不是 Clean，也不是没有诊断信息的 Dirty}
\end{table}

失败时若清 dirty，未落盘数据会被伪装成可淘汰；若只保留一个 dirty bit，又
无法分辨是否已有 writer。显式 Writeback/Error 让失败页保留 frame、identity
与重试责任。

dirty hard limit 在 MarkDirty 之前检查。达到上限，页故障返回失败，PTE 仍
只读。不能先开放 PTE 再发现无容量，否则用户已经能够继续修改，缓存却无法
记录这份债务。

\topic{用写保护捕获 MAP\_SHARED 修改}
x86 没有“每次 store 都调用 Kernel”的页模式。v1.16 使用 write-protect
notification：

\begin{enumerate}
  \item 完整 FileShared 页从 VFS 读入 FilePageCache；
  \item 所有 alias 指向同一 frame，PTE 初始只读；
  \item 第一次用户 store 产生 present+write+user fault；
  \item Kernel 验证 VMA、文件范围和 writable open，执行 Clean 到 Dirty；
  \item 成功后才设置该 PTE.W，并 \code{INVLPG}；
  \item CPU 重试原 store，其他 alias 立即看到共享 frame 的新字节。
\end{enumerate}

fork 不对 FileShared 加 COW，父子仍共享文件页；FilePrivate 继续走 COW，
不会进入文件写回。

sync 开始时必须重新写保护所有 writable FileShared PTE。否则 writer 把页
交给 VFS 后，用户仍能经 writable PTE 修改，而缓存最终被标为 Clean。重新
保护让“写回快照之后的新 store”再次 fault 并形成下一轮 Dirty。

\topic{从页身份到 ATA FLUSH 稳定边界}\par
\begin{lstlisting}
Protect shared writable PTEs
  -> Dirty/Error page Writeback
  -> UserFileBackingManager::WritePage
  -> Vfs::WriteAt(file offset; no fd offset change)
  -> rootfs transaction and BlockCache
  -> Vfs::Sync
  -> asynchronous ATA FLUSH
  -> IRQ14 Succeeded
\end{lstlisting}

文件页 identity 不能使用 fd：fd 可以关闭、duplicate 和复用，映射仍继续存在。
当前 key 使用 mount/superblock identity、inode identifier 与 inode
generation。

rootfs 有两种 generation。transaction generation 描述盘面提交次数，应随
普通写增长；VFS superblock generation 描述一次 mount 的身份，必须保持
稳定。若把前者每次复制给后者，同一 inode 会在一次写前后产生两个 cache key，
shared alias 便不再共享。v1.16 的 rootfs 集成测试专门冻结这条边界。

FLUSH 完成只说明当前设备模型报告先前写入已越过其 cache 边界；v1.17 还需
定义数据、metadata、journal descriptor 和 commit record 的顺序。没有
journal 时，v1.16 仍只能在 Dirty superblock 后拒绝不完整事务，不能自动恢复。

\topic{日志为何必须低频且带两种时间}
ATA status、PIO word、cache lookup、PIT tick 都是热路径。逐次打印会让串口
成为主要工作负载并改变 IRQ 竞争。项目只在初始化、请求提交、唯一完成、
shared writeback 验证和终态统计输出 marker。

\begin{lstlisting}
ATA_IRQ14_READY
ATA_REQUEST_CAPACITY=0x40
FLUSH_SUBMIT_REQUEST
FLUSH_SUBMIT_TIME_NS
COMPLETION_RESULT
COMPLETION_TIME_NS
OTHER_THREAD_PROGRESS
FILE_SHARED_WRITEBACK_VERIFIED
\end{lstlisting}

来宾纳秒来自 PIT，用于 deadline 和顺序；QEMU runner 添加的
\code{[QEMU][T+...ms]} 是宿主收到串口行的单调时间，只用于定位慢阶段。
两种时钟起点和调度环境不同，不能比较绝对值。

\topic{测试怎样覆盖状态而不是只看最终成功}\par
\begin{table}[htbp]
\centering
\small
\begin{osgridtabularx}{\textwidth}{l X}
证据层 & v1.16 覆盖\\ \hline
单元 & 参数、FIFO、单飞、容量、取消、完成、超时、Reap、dirty limit、Error 重试\\ \hline
集成 & completion/timeout 后下一请求、PIC cascade、rootfs identity 稳定\\ \hline
随机 & 两个固定种子模型各 100000 步，逐步对照请求与页状态参考模型\\ \hline
故障 & ATA busy/error/timeout、重复解析、writer 失败与重试\\ \hline
QEMU & 三档 RAM、真实 IRQ14、shared alias、落盘读回、private 隔离与 Thread 前进\\ \hline
\end{osgridtabularx}
\caption{最终出现一行 READY 不能代替中间状态和失败证据}
\end{table}

多个 64 GiB TCG QEMU 并发会竞争宿主 CPU，交互式逐键协议可能超过有限进度
门限。宿主逻辑测试可以并行；交互式、persistence 和 capacity QEMU 在发布时
串行复验，并继续保留有限 timeout。无限等待不能区分“机器慢”和真实死锁。

这次串行复验还抓到一个跨阶段竞态：父 Shell 已发布 foreground PGID 时，
child 可能尚未运行到恢复默认 Ctrl-C/Z disposition，于是控制信号调用从
Shell 继承的 handler。正确修复不是让 QMP 多等一会，而是父进程在 fork 前
屏蔽两种信号；child 先恢复默认 action、加入 PGID，再恢复继承前 mask。
窗口内到达的信号只进入 pending，解除 mask 时才按 child 的默认动作交付。

\topic{本阶段没有宣称什么}\par
\begin{itemize}
  \item ATA 仍是 primary master、LBA28、PIO、单飞，没有 DMA/AHCI/NVMe；
  \item writable shared 暂只接受完整页；
  \item 只有全局显式 sync，没有 \code{msync} 和后台 flusher；
  \item Error 页持续占用 cache 并产生回压，不会静默丢弃；
  \item rootfs 尚无 journal/replay；v1.17 才加入 credits、commit 与断电矩阵；
  \item 单 BSP 下完成状态正确，不宣称已经处理跨 CPU completion 与 cache coherence。
\end{itemize}

这些边界把学习顺序保持清晰：v1.16 先建立“请求只能完成一次、脏页责任不会
消失、写回有真实设备稳定边界”，v1.17 再讨论多块元数据事务怎样在断电后只
恢复为完整旧状态或完整新状态。
